\documentclass{article}
\usepackage{graphicx} 

\setlength{\parindent}{15pt}
\setlength{\parskip}{5pt}

\usepackage[margin=1.5in]{geometry}
\usepackage{amssymb}
\usepackage{amsmath}
\usepackage{mathtools}
\usepackage{xcolor}
\usepackage{cancel}
\usepackage{centernot}
\usepackage{nicefrac}
\usepackage{bm}
\usepackage{indentfirst}
\usepackage{float}
\usepackage[round]{natbib} 
\usepackage[hidelinks]{hyperref}
\usepackage{enumerate}
\usepackage{titling}
\usepackage{enumitem}
\usepackage{parskip}
\usepackage{subcaption}

\setlist[itemize]{label=\textendash}

\newenvironment{answer}{\par\vspace{0.8em}\textbf{A:}}{\par}
\newcommand{\red}[1]{\textcolor{red}{#1}}
\newcommand{\blue}[1]{\textcolor{blue}{#1}}
\newcommand{\fa}{\text{ for all }} 
\newcommand{\ex}{\text{ there exists }}
\newcommand{\st}{\text{ such that }}
\newcommand{\for}{\text{ for }}
\newcommand{\andt}{\text{ and }}
\newcommand{\ort}{\text{ or }}
\newcommand{\ift}{\text{ if }}
\newcommand{\where}{\text{ where }}
\newcommand{\but}{\text{ but }}
\newcommand{\nat}{\mathbb{N}}
\newcommand{\rat}{\mathbb{Q}}
\newcommand{\com}{\mathbb{C}}
\newcommand{\real}{\mathbb{R}}
\newcommand{\norm}[1]{\lVert#1\rVert}
\newcommand{\abs}[1]{\lvert#1\rvert}
\newcommand{\corr}{\rightrightarrows}
\newcommand{\ul}{\underline}
\newcommand{\bft}{\textbf}
\newcommand{\itt}{\textit}
\newcommand{\E}{\mathbb{E}}
\newcommand{\Prob}{\operatorname{P}}
\newcommand{\Var}{\operatorname{Var}}
\newcommand{\Cov}{\operatorname{Cov}}
\newcommand{\Unif}{\operatorname{Unif}}
\newcommand{\Bias}{\operatorname{Bias}}
\newcommand{\N}{\operatorname{N}}
\newcommand{\MSE}{\operatorname{MSE}}
\newcommand{\iffs}{\Leftrightarrow}
\newcommand{\prefeq}{\trianglerighteq}
\newcommand{\pref}{\triangleright}
\newcommand{\indiff}{\mathrel{\triangleright\triangleleft}}

\title{Next steps and future analysis}
\date{\today}

\begin{document}

\maketitle

\subsection*{Next steps}

\begin{itemize} 
    \item Start drafting; sequence findings; make slides
    \item Re-read literature for core findings about age/beliefs/redistribution to position Prolific survey and the broader project; edit lit review
    \item Investigate drivers of the belief shift
    \begin{itemize}
        \item Timing: old moved around 1990, young moved strongly around 2020--2024.
        \item State/local U.S. data to test causal effect of age inequality in economic conditions and political representation.
        \item Text data (Congressional speeches, Manifesto project) to test supply of attribution/zero-sum rich/race political narratives. These beliefs fit the villain-victim narrative framework.
    \end{itemize}
    \item Draft/pre-register 2nd Prolific wave
    \begin{itemize}
        \item An information provision or wording experiment generating exogenous belief variation. See if/how preferences move.
        \item A more comprehensive questionnaire pairing existing questions with Haidt's values, govt. trust, religiosity, and race-related beliefs.
        \item In-depth measurement and differentiation of zero-sum/attribution beliefs.
    \end{itemize}
\end{itemize}

% \subsection*{Drivers of beliefs}
%  \blue{(running list of possible mechanisms/tests)}
%  \begin{itemize}
%     \item Stagnating within-generation income growth, especially for lower-income/less-educated?
%     \item Declining absolute mobility
%     \item Racial makeup of younger cohorts
%     \item Political/institutional trajectory, specific policy choices. Increasing power of the old in U.S. politics? 
%     \item Media, social/information environment, supply of narratives. Amplifier more than root cause?
%     \item Homeownership and assets
%     \item Erosion of social insurance and secularization
%     % \item Immigration slowdown
%     % \item Self-reported belief updating questions
% \end{itemize}


\subsection*{Prolific 2nd wave}

Running list of questions that could be useful in a second wave.

Must-have:
\begin{itemize}
    \item Priming or wording or information treatment
    \begin{itemize}
        \item To shift beliefs about the economic game/zero-sum. Even to generate income/path uncertainty?
    \end{itemize}
    \item Decomposition of zero-sum wealth/rich/competition
    \begin{itemize}
        \item Fixed-pie vs. scarce positions/resources vs. ...
    \end{itemize}
    \item Distinguish between zero-sum and deservingness/attribution
    \begin{itemize}
        \item Zero-sum: how the income distribution changes in response to someone's personal effort or gains. Deservingness: whether those who benefit from an income draw or program are seen as responsible for their situation.
    \end{itemize}
    \item Why someone becomes wealthy
    \begin{itemize}
        \item (i) effort, investment, risk-taking, innovation (ii) structural advantages, rich family (iii) luck. 
        \item Structural advantages questions seem highly significant in other surveys, especially if concrete examples are mentioned. Run a battery of advantages: education, race, wealth, connections, etc.
    \end{itemize}
    \item Why someone becomes poor
    \begin{itemize}
        \item Because they are lazy or unlucky?
    \end{itemize}
    \item What income should track 
    \begin{itemize}
        \item (i) effort, time, responsibility, seniority (ii) skills, intelligence, credentials, investment (iii) performance.
    \end{itemize}
    \item Zero-sum race/disadvantaged  
    \begin{itemize}
        \item More about minorities or poorer groups or forced equity. DEI is attached to gender, sexual identity, etc.
    \end{itemize}
    \item Racism and discrimination
    \item Perceived conflict/antagonism between rich and poor
    \item Deeper questions about whether/why one's beliefs changed
    \begin{itemize}
        \item Asking about zero-sum, incentive, and mobility beliefs, not just redistributive preferences.
    \end{itemize}
    \item Reasons for supporting/opposing redistribution in their own words
    \begin{itemize}
        \item Easier to extract insights from free-response questions with LLMs.
    \end{itemize}
    \item Missing values/beliefs questions
    \begin{itemize}
        \item Religiosity (importance in life, attendance), social class (self-reported), opinions on immigration, trust/confidence in government (including capacity/capability), generalized social trust, political/civic involvement.
    \end{itemize}
    \item Missing economic/demographic information
    \begin{itemize}
        \item Health insurance coverage, homeownership, wealth and previous income, retirement status and benefits.
    \end{itemize}
    \item Remove policy from question stem as much as possible
    \begin{itemize}
        \item Especially for gapmotivates and taxdemotivates.
    \end{itemize}
    \item MFT/Haidt moral values
    \begin{itemize}
        \item Care/compassion, proportionality, equality, obligation/authority, loyalty, purity/sanctity
    \end{itemize}
\end{itemize}

Nice to have:
\begin{itemize}
    \item Precise measurement of generosity
    \begin{itemize}
        \item Separate it from universalism, scale with income if numeric.
    \end{itemize}
    \item Decomposition of mobility beliefs
    \begin{itemize}
        \item Absolute vs. relative mobility.  Personal mobility: I or people like me can/can't move up. Structural mobility: origins do/don't constrain outcomes. Meritocratic mobility: moving up is/isn't based on talent or providing value.
    \end{itemize}
    \item Separate individual vs aggregate incentive effects
    \begin{itemize}
        \item The distinction seems to matter in Stantcheva 2021.
    \end{itemize}
    \item More-precise questions about financial expectations, income/wealth variance, and financial struggles
    \begin{itemize}
        \item Some other papers and longitudinal surveys find financial situation significant when measured certain ways.
    \end{itemize}
    \item Separate myopia (personality trait) from need for money now (financial situation)
    \item Ask about policy variations
    \begin{itemize}
        \item To single out preference for ex-ante insurance vs ex-post redistribution. 
    \end{itemize}
    \item Tax incidence
    \begin{itemize}
        \item Who, in the end, benefits from redistributive taxes or higher taxes?
    \end{itemize}
    \item Post-materialism 
    \begin{itemize}
        \item climaterisk could be capturing a latent worldview/ideology (prioritizing goals other than economic growth)
    \end{itemize}
    \item Visual tax proposal without experimental variation
    \begin{itemize}
        \item Our current experimental variation introduces noise in the regressions
    \end{itemize}
\end{itemize}





\end{document}
